% Options for packages loaded elsewhere
% Options for packages loaded elsewhere
\PassOptionsToPackage{unicode}{hyperref}
\PassOptionsToPackage{hyphens}{url}
\PassOptionsToPackage{dvipsnames,svgnames,x11names}{xcolor}
%
\documentclass[
  letterpaper,
  DIV=11,
  numbers=noendperiod]{scrreprt}
\usepackage{xcolor}
\usepackage{amsmath,amssymb}
\setcounter{secnumdepth}{5}
\usepackage{iftex}
\ifPDFTeX
  \usepackage[T1]{fontenc}
  \usepackage[utf8]{inputenc}
  \usepackage{textcomp} % provide euro and other symbols
\else % if luatex or xetex
  \usepackage{unicode-math} % this also loads fontspec
  \defaultfontfeatures{Scale=MatchLowercase}
  \defaultfontfeatures[\rmfamily]{Ligatures=TeX,Scale=1}
\fi
\usepackage{lmodern}
\ifPDFTeX\else
  % xetex/luatex font selection
\fi
% Use upquote if available, for straight quotes in verbatim environments
\IfFileExists{upquote.sty}{\usepackage{upquote}}{}
\IfFileExists{microtype.sty}{% use microtype if available
  \usepackage[]{microtype}
  \UseMicrotypeSet[protrusion]{basicmath} % disable protrusion for tt fonts
}{}
\makeatletter
\@ifundefined{KOMAClassName}{% if non-KOMA class
  \IfFileExists{parskip.sty}{%
    \usepackage{parskip}
  }{% else
    \setlength{\parindent}{0pt}
    \setlength{\parskip}{6pt plus 2pt minus 1pt}}
}{% if KOMA class
  \KOMAoptions{parskip=half}}
\makeatother
% Make \paragraph and \subparagraph free-standing
\makeatletter
\ifx\paragraph\undefined\else
  \let\oldparagraph\paragraph
  \renewcommand{\paragraph}{
    \@ifstar
      \xxxParagraphStar
      \xxxParagraphNoStar
  }
  \newcommand{\xxxParagraphStar}[1]{\oldparagraph*{#1}\mbox{}}
  \newcommand{\xxxParagraphNoStar}[1]{\oldparagraph{#1}\mbox{}}
\fi
\ifx\subparagraph\undefined\else
  \let\oldsubparagraph\subparagraph
  \renewcommand{\subparagraph}{
    \@ifstar
      \xxxSubParagraphStar
      \xxxSubParagraphNoStar
  }
  \newcommand{\xxxSubParagraphStar}[1]{\oldsubparagraph*{#1}\mbox{}}
  \newcommand{\xxxSubParagraphNoStar}[1]{\oldsubparagraph{#1}\mbox{}}
\fi
\makeatother

\usepackage{color}
\usepackage{fancyvrb}
\newcommand{\VerbBar}{|}
\newcommand{\VERB}{\Verb[commandchars=\\\{\}]}
\DefineVerbatimEnvironment{Highlighting}{Verbatim}{commandchars=\\\{\}}
% Add ',fontsize=\small' for more characters per line
\usepackage{framed}
\definecolor{shadecolor}{RGB}{241,243,245}
\newenvironment{Shaded}{\begin{snugshade}}{\end{snugshade}}
\newcommand{\AlertTok}[1]{\textcolor[rgb]{0.68,0.00,0.00}{#1}}
\newcommand{\AnnotationTok}[1]{\textcolor[rgb]{0.37,0.37,0.37}{#1}}
\newcommand{\AttributeTok}[1]{\textcolor[rgb]{0.40,0.45,0.13}{#1}}
\newcommand{\BaseNTok}[1]{\textcolor[rgb]{0.68,0.00,0.00}{#1}}
\newcommand{\BuiltInTok}[1]{\textcolor[rgb]{0.00,0.23,0.31}{#1}}
\newcommand{\CharTok}[1]{\textcolor[rgb]{0.13,0.47,0.30}{#1}}
\newcommand{\CommentTok}[1]{\textcolor[rgb]{0.37,0.37,0.37}{#1}}
\newcommand{\CommentVarTok}[1]{\textcolor[rgb]{0.37,0.37,0.37}{\textit{#1}}}
\newcommand{\ConstantTok}[1]{\textcolor[rgb]{0.56,0.35,0.01}{#1}}
\newcommand{\ControlFlowTok}[1]{\textcolor[rgb]{0.00,0.23,0.31}{\textbf{#1}}}
\newcommand{\DataTypeTok}[1]{\textcolor[rgb]{0.68,0.00,0.00}{#1}}
\newcommand{\DecValTok}[1]{\textcolor[rgb]{0.68,0.00,0.00}{#1}}
\newcommand{\DocumentationTok}[1]{\textcolor[rgb]{0.37,0.37,0.37}{\textit{#1}}}
\newcommand{\ErrorTok}[1]{\textcolor[rgb]{0.68,0.00,0.00}{#1}}
\newcommand{\ExtensionTok}[1]{\textcolor[rgb]{0.00,0.23,0.31}{#1}}
\newcommand{\FloatTok}[1]{\textcolor[rgb]{0.68,0.00,0.00}{#1}}
\newcommand{\FunctionTok}[1]{\textcolor[rgb]{0.28,0.35,0.67}{#1}}
\newcommand{\ImportTok}[1]{\textcolor[rgb]{0.00,0.46,0.62}{#1}}
\newcommand{\InformationTok}[1]{\textcolor[rgb]{0.37,0.37,0.37}{#1}}
\newcommand{\KeywordTok}[1]{\textcolor[rgb]{0.00,0.23,0.31}{\textbf{#1}}}
\newcommand{\NormalTok}[1]{\textcolor[rgb]{0.00,0.23,0.31}{#1}}
\newcommand{\OperatorTok}[1]{\textcolor[rgb]{0.37,0.37,0.37}{#1}}
\newcommand{\OtherTok}[1]{\textcolor[rgb]{0.00,0.23,0.31}{#1}}
\newcommand{\PreprocessorTok}[1]{\textcolor[rgb]{0.68,0.00,0.00}{#1}}
\newcommand{\RegionMarkerTok}[1]{\textcolor[rgb]{0.00,0.23,0.31}{#1}}
\newcommand{\SpecialCharTok}[1]{\textcolor[rgb]{0.37,0.37,0.37}{#1}}
\newcommand{\SpecialStringTok}[1]{\textcolor[rgb]{0.13,0.47,0.30}{#1}}
\newcommand{\StringTok}[1]{\textcolor[rgb]{0.13,0.47,0.30}{#1}}
\newcommand{\VariableTok}[1]{\textcolor[rgb]{0.07,0.07,0.07}{#1}}
\newcommand{\VerbatimStringTok}[1]{\textcolor[rgb]{0.13,0.47,0.30}{#1}}
\newcommand{\WarningTok}[1]{\textcolor[rgb]{0.37,0.37,0.37}{\textit{#1}}}

\usepackage{longtable,booktabs,array}
\usepackage{calc} % for calculating minipage widths
% Correct order of tables after \paragraph or \subparagraph
\usepackage{etoolbox}
\makeatletter
\patchcmd\longtable{\par}{\if@noskipsec\mbox{}\fi\par}{}{}
\makeatother
% Allow footnotes in longtable head/foot
\IfFileExists{footnotehyper.sty}{\usepackage{footnotehyper}}{\usepackage{footnote}}
\makesavenoteenv{longtable}
\usepackage{graphicx}
\makeatletter
\newsavebox\pandoc@box
\newcommand*\pandocbounded[1]{% scales image to fit in text height/width
  \sbox\pandoc@box{#1}%
  \Gscale@div\@tempa{\textheight}{\dimexpr\ht\pandoc@box+\dp\pandoc@box\relax}%
  \Gscale@div\@tempb{\linewidth}{\wd\pandoc@box}%
  \ifdim\@tempb\p@<\@tempa\p@\let\@tempa\@tempb\fi% select the smaller of both
  \ifdim\@tempa\p@<\p@\scalebox{\@tempa}{\usebox\pandoc@box}%
  \else\usebox{\pandoc@box}%
  \fi%
}
% Set default figure placement to htbp
\def\fps@figure{htbp}
\makeatother





\setlength{\emergencystretch}{3em} % prevent overfull lines

\providecommand{\tightlist}{%
  \setlength{\itemsep}{0pt}\setlength{\parskip}{0pt}}



 


\KOMAoption{captions}{tableheading}
\makeatletter
\@ifpackageloaded{bookmark}{}{\usepackage{bookmark}}
\makeatother
\makeatletter
\@ifpackageloaded{caption}{}{\usepackage{caption}}
\AtBeginDocument{%
\ifdefined\contentsname
  \renewcommand*\contentsname{Table of contents}
\else
  \newcommand\contentsname{Table of contents}
\fi
\ifdefined\listfigurename
  \renewcommand*\listfigurename{List of Figures}
\else
  \newcommand\listfigurename{List of Figures}
\fi
\ifdefined\listtablename
  \renewcommand*\listtablename{List of Tables}
\else
  \newcommand\listtablename{List of Tables}
\fi
\ifdefined\figurename
  \renewcommand*\figurename{Figure}
\else
  \newcommand\figurename{Figure}
\fi
\ifdefined\tablename
  \renewcommand*\tablename{Table}
\else
  \newcommand\tablename{Table}
\fi
}
\@ifpackageloaded{float}{}{\usepackage{float}}
\floatstyle{ruled}
\@ifundefined{c@chapter}{\newfloat{codelisting}{h}{lop}}{\newfloat{codelisting}{h}{lop}[chapter]}
\floatname{codelisting}{Listing}
\newcommand*\listoflistings{\listof{codelisting}{List of Listings}}
\makeatother
\makeatletter
\makeatother
\makeatletter
\@ifpackageloaded{caption}{}{\usepackage{caption}}
\@ifpackageloaded{subcaption}{}{\usepackage{subcaption}}
\makeatother
\usepackage{bookmark}
\IfFileExists{xurl.sty}{\usepackage{xurl}}{} % add URL line breaks if available
\urlstyle{same}
\hypersetup{
  pdftitle={diffusion algorithms for bacteria},
  pdfauthor={Norah Jones},
  colorlinks=true,
  linkcolor={blue},
  filecolor={Maroon},
  citecolor={Blue},
  urlcolor={Blue},
  pdfcreator={LaTeX via pandoc}}


\title{diffusion algorithms for bacteria}
\author{Norah Jones}
\date{2026-08-20}
\begin{document}
\maketitle

\renewcommand*\contentsname{Table of contents}
{
\hypersetup{linkcolor=}
\setcounter{tocdepth}{2}
\tableofcontents
}

\bookmarksetup{startatroot}

\chapter*{Preface}\label{preface}
\addcontentsline{toc}{chapter}{Preface}

\markboth{Preface}{Preface}

This is a Quarto book.

To learn more about Quarto books visit
\url{https://quarto.org/docs/books}.

\bookmarksetup{startatroot}

\chapter{Explaining the ADI Method}\label{explaining-the-adi-method}

The Alternating Direction Implicit (ADI) method solves the 2‑D diffusion
equation with optional sink terms. It splits each time step into two
half‑steps: an \textbf{x‑sweep} that treats the x‑derivative implicitly
and the y‑derivative explicitly, followed by a \textbf{y‑sweep} that
does the opposite. This turns a large 2‑D system into many small
tridiagonal systems that can be solved efficiently with the Thomas
algorithm.

The module keeps all working arrays and coefficients at module scope so
that they are allocated only once and reused across many calls to
\texttt{ADI}.

\section{Imports and module‑level
state}\label{imports-and-modulelevel-state}

We import the Thomas algorithm solver and a helper that pre‑allocates
all the arrays we need.

\phantomsection\label{adi-imports}%
\begin{Shaded}
\begin{Highlighting}[]
\ImportTok{import}\NormalTok{ \{ thomasAlgorithm \} }\ImportTok{from} \StringTok{"./thomasAlgorithm.js"}\OperatorTok{;}
\ImportTok{import}\NormalTok{ \{ initADIArrays \} }\ImportTok{from} \StringTok{"./initArrays.js"}\OperatorTok{;}
\end{Highlighting}
\end{Shaded}

The following variables hold the grid dimensions, the tridiagonal
coefficients for both sweeps, the intermediate concentration field, and
the pre‑computed constants \texttt{alpha}, \texttt{halfDeltaT},
\texttt{scaledSources}, and \texttt{gamma} (which incorporates the sink
term).

\phantomsection\label{adi-variables}%
\begin{Shaded}
\begin{Highlighting}[]
\KeywordTok{let}\NormalTok{ WIDTH}\OperatorTok{;}
\KeywordTok{let}\NormalTok{ HEIGHT}\OperatorTok{;}
\KeywordTok{let}\NormalTok{ modifiedUpperDiagonal1}\OperatorTok{,}\NormalTok{ modifiedRightHandSide1}\OperatorTok{,}\NormalTok{ solution1}\OperatorTok{;}
\KeywordTok{let}\NormalTok{ modifiedUpperDiagonal2}\OperatorTok{,}\NormalTok{ modifiedRightHandSide2}\OperatorTok{,}\NormalTok{ solution2}\OperatorTok{;}
\KeywordTok{let}\NormalTok{ intermediateConcentration}\OperatorTok{;}
\KeywordTok{let}\NormalTok{ a1}\OperatorTok{,}\NormalTok{ b1}\OperatorTok{,}\NormalTok{ c1}\OperatorTok{,}\NormalTok{ d1}\OperatorTok{;}
\KeywordTok{let}\NormalTok{ a2}\OperatorTok{,}\NormalTok{ b2}\OperatorTok{,}\NormalTok{ c2}\OperatorTok{,}\NormalTok{ d2}\OperatorTok{;}
\KeywordTok{let}\NormalTok{ alpha}\OperatorTok{,}\NormalTok{ halfDeltaT}\OperatorTok{,}\NormalTok{ scaledSources}\OperatorTok{;}
\KeywordTok{let}\NormalTok{ gamma}\OperatorTok{;}
\end{Highlighting}
\end{Shaded}

\section{Initialising the ADI arrays}\label{initialising-the-adi-arrays}

\texttt{setADIProperties} must be called once before any simulation. It
stores the grid size and calls \texttt{initADIArrays} to allocate all
working storage and compute the constant
\texttt{alpha\ =\ D\ Δt\ /\ (2\ Δx²)}.

\phantomsection\label{adi-set-properties}%
\begin{Shaded}
\begin{Highlighting}[]
\ImportTok{export} \KeywordTok{const}\NormalTok{ setADIProperties }\OperatorTok{=}\NormalTok{ (}
\NormalTok{    width}\OperatorTok{,}
\NormalTok{    height}\OperatorTok{,}
\NormalTok{    diffusionCoefficient}\OperatorTok{,}
\NormalTok{    deltaX}\OperatorTok{,}
\NormalTok{    deltaT}\OperatorTok{,}
\NormalTok{) }\KeywordTok{=\textgreater{}}\NormalTok{ \{}
\NormalTok{    WIDTH }\OperatorTok{=}\NormalTok{ width}\OperatorTok{;}
\NormalTok{    HEIGHT }\OperatorTok{=}\NormalTok{ height}\OperatorTok{;}
\NormalTok{    (\{}
\NormalTok{        modifiedUpperDiagonal1}\OperatorTok{,}
\NormalTok{        modifiedRightHandSide1}\OperatorTok{,}
\NormalTok{        solution1}\OperatorTok{,}
\NormalTok{        modifiedUpperDiagonal2}\OperatorTok{,}
\NormalTok{        modifiedRightHandSide2}\OperatorTok{,}
\NormalTok{        solution2}\OperatorTok{,}
\NormalTok{        intermediateConcentration}\OperatorTok{,}
\NormalTok{        a1}\OperatorTok{,}
\NormalTok{        b1}\OperatorTok{,}
\NormalTok{        c1}\OperatorTok{,}
\NormalTok{        d1}\OperatorTok{,}
\NormalTok{        a2}\OperatorTok{,}
\NormalTok{        b2}\OperatorTok{,}
\NormalTok{        c2}\OperatorTok{,}
\NormalTok{        d2}\OperatorTok{,}
\NormalTok{        alpha}\OperatorTok{,}
\NormalTok{        halfDeltaT}\OperatorTok{,}
\NormalTok{        scaledSources}\OperatorTok{,}
\NormalTok{        gamma}\OperatorTok{,}
\NormalTok{    \} }\OperatorTok{=} \FunctionTok{initADIArrays}\NormalTok{(WIDTH}\OperatorTok{,}\NormalTok{ HEIGHT}\OperatorTok{,}\NormalTok{ diffusionCoefficient}\OperatorTok{,}\NormalTok{ deltaX}\OperatorTok{,}\NormalTok{ deltaT))}\OperatorTok{;}
\NormalTok{\}}\OperatorTok{;}
\end{Highlighting}
\end{Shaded}

\section{Updating sinks and sources}\label{updating-sinks-and-sources}

Before each call to \texttt{ADI} (or whenever the source/sink arrays
change), \texttt{updateSinksAndSources} pre‑scales the source and sink
terms so that they can be added directly to the right‑hand side during
the sweeps.

\phantomsection\label{adi-update-sinks-sources}%
\begin{Shaded}
\begin{Highlighting}[]
\ImportTok{export} \KeywordTok{const}\NormalTok{ updateSinksAndSources }\OperatorTok{=}\NormalTok{ (sinks}\OperatorTok{,}\NormalTok{ sources) }\KeywordTok{=\textgreater{}}\NormalTok{ \{}
    \ControlFlowTok{for}\NormalTok{ (}\KeywordTok{let}\NormalTok{ i }\OperatorTok{=} \DecValTok{0}\OperatorTok{;}\NormalTok{ i }\OperatorTok{\textless{}}\NormalTok{ WIDTH }\OperatorTok{*}\NormalTok{ HEIGHT}\OperatorTok{;}\NormalTok{ i}\OperatorTok{++}\NormalTok{) \{}
\NormalTok{        gamma[i] }\OperatorTok{=}\NormalTok{ (sinks[i] }\OperatorTok{*}\NormalTok{ halfDeltaT) }\OperatorTok{/} \DecValTok{2}\OperatorTok{;}
\NormalTok{        scaledSources[i] }\OperatorTok{=}\NormalTok{ sources[i] }\OperatorTok{*}\NormalTok{ halfDeltaT}\OperatorTok{;}
\NormalTok{    \}}
\NormalTok{\}}
\end{Highlighting}
\end{Shaded}

\section{Updating the main diagonal for each
sweep}\label{updating-the-main-diagonal-for-each-sweep}

Because the sink term \texttt{gamma} varies in space, the main diagonal
of the tridiagonal system must be rebuilt for every row (x‑sweep) or
column (y‑sweep). Neumann (zero‑flux) boundary conditions are enforced
by adjusting the first and last diagonal entries.

\phantomsection\label{adi-update-diagonal-x}%
\begin{Shaded}
\begin{Highlighting}[]
\KeywordTok{const}\NormalTok{ updateMainDiagonalXstep }\OperatorTok{=}\NormalTok{ (yCoord) }\KeywordTok{=\textgreater{}}\NormalTok{ \{}
    \CommentTok{//update b1}
    \ControlFlowTok{for}\NormalTok{ (}\KeywordTok{let}\NormalTok{ i }\OperatorTok{=} \DecValTok{0}\OperatorTok{;}\NormalTok{ i }\OperatorTok{\textless{}}\NormalTok{ WIDTH}\OperatorTok{;}\NormalTok{ i}\OperatorTok{++}\NormalTok{) \{}
        \KeywordTok{const}\NormalTok{ idx }\OperatorTok{=}\NormalTok{ yCoord }\OperatorTok{*}\NormalTok{ WIDTH }\OperatorTok{+}\NormalTok{ i}\OperatorTok{;}
\NormalTok{        b1[i] }\OperatorTok{=} \DecValTok{1} \OperatorTok{+} \DecValTok{2} \OperatorTok{*}\NormalTok{ alpha }\OperatorTok{+}\NormalTok{ gamma[idx]}\OperatorTok{;}
\NormalTok{    \}}
    \CommentTok{//boundary conditions}
\NormalTok{    b1[}\DecValTok{0}\NormalTok{] }\OperatorTok{=} \DecValTok{1} \OperatorTok{+}\NormalTok{ alpha }\OperatorTok{+}\NormalTok{ gamma[yCoord }\OperatorTok{*}\NormalTok{ WIDTH }\OperatorTok{+} \DecValTok{0}\NormalTok{]}\OperatorTok{;}
\NormalTok{    b1[WIDTH }\OperatorTok{{-}} \DecValTok{1}\NormalTok{] }\OperatorTok{=} \DecValTok{1} \OperatorTok{+}\NormalTok{ alpha }\OperatorTok{+}\NormalTok{ gamma[yCoord }\OperatorTok{*}\NormalTok{ WIDTH }\OperatorTok{+}\NormalTok{ (WIDTH }\OperatorTok{{-}} \DecValTok{1}\NormalTok{)]}\OperatorTok{;}
\NormalTok{\}}
\end{Highlighting}
\end{Shaded}

\phantomsection\label{adi-update-diagonal-y}%
\begin{Shaded}
\begin{Highlighting}[]
\KeywordTok{const}\NormalTok{ updateMainDiagonalYstep }\OperatorTok{=}\NormalTok{ (xCoord) }\KeywordTok{=\textgreater{}}\NormalTok{ \{}
    \CommentTok{//update b2}
    \ControlFlowTok{for}\NormalTok{ (}\KeywordTok{let}\NormalTok{ j }\OperatorTok{=} \DecValTok{0}\OperatorTok{;}\NormalTok{ j }\OperatorTok{\textless{}}\NormalTok{ HEIGHT}\OperatorTok{;}\NormalTok{ j}\OperatorTok{++}\NormalTok{) \{}
        \KeywordTok{const}\NormalTok{ idx }\OperatorTok{=}\NormalTok{ j }\OperatorTok{*}\NormalTok{ WIDTH }\OperatorTok{+}\NormalTok{ xCoord}\OperatorTok{;}
\NormalTok{        b2[j] }\OperatorTok{=} \DecValTok{1} \OperatorTok{+} \DecValTok{2} \OperatorTok{*}\NormalTok{ alpha }\OperatorTok{+}\NormalTok{ gamma[idx]}\OperatorTok{;}
\NormalTok{    \}}
    \CommentTok{//boundary conditions}
\NormalTok{    b2[}\DecValTok{0}\NormalTok{] }\OperatorTok{=} \DecValTok{1} \OperatorTok{+}\NormalTok{ alpha }\OperatorTok{+}\NormalTok{ gamma[}\DecValTok{0} \OperatorTok{*}\NormalTok{ WIDTH }\OperatorTok{+}\NormalTok{ xCoord]}\OperatorTok{;}
\NormalTok{    b2[HEIGHT }\OperatorTok{{-}} \DecValTok{1}\NormalTok{] }\OperatorTok{=} \DecValTok{1} \OperatorTok{+}\NormalTok{ alpha }\OperatorTok{+}\NormalTok{ gamma[(HEIGHT }\OperatorTok{{-}} \DecValTok{1}\NormalTok{) }\OperatorTok{*}\NormalTok{ WIDTH }\OperatorTok{+}\NormalTok{ xCoord]}\OperatorTok{;}
\NormalTok{\}}
\end{Highlighting}
\end{Shaded}

\section{The ADI time‑stepping loop}\label{the-adi-timestepping-loop}

The main \texttt{ADI} function advances the concentration field by
\texttt{totalNumberOfIterations} half‑time steps. Each full iteration
consists of an x‑sweep followed by a y‑sweep. Clamped indexing
(\texttt{jBelow\ =\ max(j-1,0)}, etc.) implements Neumann boundary
conditions for the explicit part.

If a negative concentration appears and \texttt{allowNegativeValues} is
\texttt{false}, the function logs a warning and returns \texttt{null}.

\phantomsection\label{adi-main-function}%
\begin{Shaded}
\begin{Highlighting}[]
\ImportTok{export} \KeywordTok{const}\NormalTok{ ADI }\OperatorTok{=}\NormalTok{ (}
\NormalTok{    concentrationData}\OperatorTok{,}
\NormalTok{    totalNumberOfIterations}\OperatorTok{,}
\NormalTok{    allowNegativeValues }\OperatorTok{=} \KeywordTok{false}
\NormalTok{) }\KeywordTok{=\textgreater{}}\NormalTok{ \{}
    \KeywordTok{let}\NormalTok{ reachedNegativeValue }\OperatorTok{=} \KeywordTok{false}\OperatorTok{;}

    \KeywordTok{const}\NormalTok{ currentConcentrationData }\OperatorTok{=}\NormalTok{ concentrationData}\OperatorTok{;}

    \ControlFlowTok{for}\NormalTok{ (}\KeywordTok{let}\NormalTok{ iteration }\OperatorTok{=} \DecValTok{0}\OperatorTok{;}\NormalTok{ iteration }\OperatorTok{\textless{}}\NormalTok{ totalNumberOfIterations}\OperatorTok{;}\NormalTok{ iteration}\OperatorTok{++}\NormalTok{) \{}
        \CommentTok{/////////////{-}{-}{-}{-}{-}  FIRST HALF{-}STEP (X{-}sweep)  {-}{-}{-}{-}{-}/////////////}
        \CommentTok{// Solve implicitly in x, explicitly in y.}
        \CommentTok{// Clamped indexing implements Neumann (zero{-}flux) BCs in y:}
        \CommentTok{//   jBelow = max(j{-}1, 0),  jAbove = min(j+1, HEIGHT{-}1)}
        \ControlFlowTok{for}\NormalTok{ (}\KeywordTok{let}\NormalTok{ j }\OperatorTok{=} \DecValTok{0}\OperatorTok{;}\NormalTok{ j }\OperatorTok{\textless{}}\NormalTok{ HEIGHT}\OperatorTok{;}\NormalTok{ j}\OperatorTok{++}\NormalTok{) \{}
            \KeywordTok{const}\NormalTok{ jBelow }\OperatorTok{=}\NormalTok{ j }\OperatorTok{\textgreater{}} \DecValTok{0} \OperatorTok{?}\NormalTok{ j }\OperatorTok{{-}} \DecValTok{1} \OperatorTok{:} \DecValTok{0}\OperatorTok{;}
            \KeywordTok{const}\NormalTok{ jAbove }\OperatorTok{=}\NormalTok{ j }\OperatorTok{\textless{}}\NormalTok{ HEIGHT }\OperatorTok{{-}} \DecValTok{1} \OperatorTok{?}\NormalTok{ j }\OperatorTok{+} \DecValTok{1} \OperatorTok{:}\NormalTok{ HEIGHT }\OperatorTok{{-}} \DecValTok{1}\OperatorTok{;}
            \KeywordTok{const}\NormalTok{ rowOffset }\OperatorTok{=}\NormalTok{ j }\OperatorTok{*}\NormalTok{ WIDTH}\OperatorTok{;}

            \ControlFlowTok{for}\NormalTok{ (}\KeywordTok{let}\NormalTok{ i }\OperatorTok{=} \DecValTok{0}\OperatorTok{;}\NormalTok{ i }\OperatorTok{\textless{}}\NormalTok{ WIDTH}\OperatorTok{;}\NormalTok{ i}\OperatorTok{++}\NormalTok{) \{}
                \KeywordTok{const}\NormalTok{ idx }\OperatorTok{=}\NormalTok{ rowOffset }\OperatorTok{+}\NormalTok{ i}\OperatorTok{;}

                \KeywordTok{const}\NormalTok{ center }\OperatorTok{=}\NormalTok{ currentConcentrationData[idx]}\OperatorTok{;}
                \KeywordTok{const}\NormalTok{ bottom }\OperatorTok{=}\NormalTok{ currentConcentrationData[jBelow }\OperatorTok{*}\NormalTok{ WIDTH }\OperatorTok{+}\NormalTok{ i]}\OperatorTok{;}
                \KeywordTok{const}\NormalTok{ top }\OperatorTok{=}\NormalTok{ currentConcentrationData[jAbove }\OperatorTok{*}\NormalTok{ WIDTH }\OperatorTok{+}\NormalTok{ i]}\OperatorTok{;}

\NormalTok{                d1[i] }\OperatorTok{=}
\NormalTok{                    alpha }\OperatorTok{*}\NormalTok{ bottom }\OperatorTok{+}
\NormalTok{                    (}\DecValTok{1} \OperatorTok{{-}} \DecValTok{2} \OperatorTok{*}\NormalTok{ alpha }\OperatorTok{{-}}\NormalTok{ gamma[idx]) }\OperatorTok{*}\NormalTok{ center }\OperatorTok{+}
\NormalTok{                    alpha }\OperatorTok{*}\NormalTok{ top }\OperatorTok{+}
\NormalTok{                    scaledSources[idx]}\OperatorTok{;}
\NormalTok{            \}}
            \FunctionTok{updateMainDiagonalXstep}\NormalTok{(j)}\OperatorTok{;}

            \FunctionTok{thomasAlgorithm}\NormalTok{(}
\NormalTok{                a1}\OperatorTok{,}
\NormalTok{                b1}\OperatorTok{,}
\NormalTok{                c1}\OperatorTok{,}
\NormalTok{                d1}\OperatorTok{,}
\NormalTok{                WIDTH}\OperatorTok{,}
\NormalTok{                modifiedUpperDiagonal1}\OperatorTok{,}
\NormalTok{                modifiedRightHandSide1}\OperatorTok{,}
\NormalTok{                solution1}
\NormalTok{            )}\OperatorTok{;}

            \ControlFlowTok{for}\NormalTok{ (}\KeywordTok{let}\NormalTok{ i }\OperatorTok{=} \DecValTok{0}\OperatorTok{;}\NormalTok{ i }\OperatorTok{\textless{}}\NormalTok{ WIDTH}\OperatorTok{;}\NormalTok{ i}\OperatorTok{++}\NormalTok{) \{}
\NormalTok{                intermediateConcentration[rowOffset }\OperatorTok{+}\NormalTok{ i] }\OperatorTok{=}\NormalTok{ solution1[i]}\OperatorTok{;}
\NormalTok{            \}}
\NormalTok{        \}}

        \CommentTok{/////////////{-}{-}{-}{-}{-}  SECOND HALF{-}STEP (Y{-}sweep)  {-}{-}{-}{-}{-}/////////////}
        \CommentTok{// Solve implicitly in y, explicitly in x.}
        \CommentTok{// Clamped indexing implements Neumann (zero{-}flux) BCs in x:}
        \CommentTok{//   iLeft = max(i{-}1, 0),  iRight = min(i+1, WIDTH{-}1)}
        \ControlFlowTok{for}\NormalTok{ (}\KeywordTok{let}\NormalTok{ i }\OperatorTok{=} \DecValTok{0}\OperatorTok{;}\NormalTok{ i }\OperatorTok{\textless{}}\NormalTok{ WIDTH}\OperatorTok{;}\NormalTok{ i}\OperatorTok{++}\NormalTok{) \{}
            \KeywordTok{const}\NormalTok{ iLeft }\OperatorTok{=}\NormalTok{ i }\OperatorTok{\textgreater{}} \DecValTok{0} \OperatorTok{?}\NormalTok{ i }\OperatorTok{{-}} \DecValTok{1} \OperatorTok{:} \DecValTok{0}\OperatorTok{;}
            \KeywordTok{const}\NormalTok{ iRight }\OperatorTok{=}\NormalTok{ i }\OperatorTok{\textless{}}\NormalTok{ WIDTH }\OperatorTok{{-}} \DecValTok{1} \OperatorTok{?}\NormalTok{ i }\OperatorTok{+} \DecValTok{1} \OperatorTok{:}\NormalTok{ WIDTH }\OperatorTok{{-}} \DecValTok{1}\OperatorTok{;}

            \ControlFlowTok{for}\NormalTok{ (}\KeywordTok{let}\NormalTok{ j }\OperatorTok{=} \DecValTok{0}\OperatorTok{;}\NormalTok{ j }\OperatorTok{\textless{}}\NormalTok{ HEIGHT}\OperatorTok{;}\NormalTok{ j}\OperatorTok{++}\NormalTok{) \{}
                \KeywordTok{const}\NormalTok{ rowOffset }\OperatorTok{=}\NormalTok{ j }\OperatorTok{*}\NormalTok{ WIDTH}\OperatorTok{;}
                \KeywordTok{const}\NormalTok{ idx }\OperatorTok{=}\NormalTok{ rowOffset }\OperatorTok{+}\NormalTok{ i}\OperatorTok{;}

                \KeywordTok{const}\NormalTok{ center }\OperatorTok{=}\NormalTok{ intermediateConcentration[idx]}\OperatorTok{;}
                \KeywordTok{const}\NormalTok{ left }\OperatorTok{=}\NormalTok{ intermediateConcentration[rowOffset }\OperatorTok{+}\NormalTok{ iLeft]}\OperatorTok{;}
                \KeywordTok{const}\NormalTok{ right }\OperatorTok{=}\NormalTok{ intermediateConcentration[rowOffset }\OperatorTok{+}\NormalTok{ iRight]}\OperatorTok{;}

\NormalTok{                d2[j] }\OperatorTok{=}
\NormalTok{                    alpha }\OperatorTok{*}\NormalTok{ left }\OperatorTok{+}
\NormalTok{                    (}\DecValTok{1} \OperatorTok{{-}} \DecValTok{2} \OperatorTok{*}\NormalTok{ alpha }\OperatorTok{{-}}\NormalTok{ gamma[idx]) }\OperatorTok{*}\NormalTok{ center }\OperatorTok{+}
\NormalTok{                    alpha }\OperatorTok{*}\NormalTok{ right }\OperatorTok{+}
\NormalTok{                    scaledSources[idx]}\OperatorTok{;}
\NormalTok{            \}}
            \FunctionTok{updateMainDiagonalYstep}\NormalTok{(i)}\OperatorTok{;}

            \FunctionTok{thomasAlgorithm}\NormalTok{(}
\NormalTok{                a2}\OperatorTok{,}
\NormalTok{                b2}\OperatorTok{,}
\NormalTok{                c2}\OperatorTok{,}
\NormalTok{                d2}\OperatorTok{,}
\NormalTok{                HEIGHT}\OperatorTok{,}
\NormalTok{                modifiedUpperDiagonal2}\OperatorTok{,}
\NormalTok{                modifiedRightHandSide2}\OperatorTok{,}
\NormalTok{                solution2}
\NormalTok{            )}\OperatorTok{;}

            \ControlFlowTok{for}\NormalTok{ (}\KeywordTok{let}\NormalTok{ j }\OperatorTok{=} \DecValTok{0}\OperatorTok{;}\NormalTok{ j }\OperatorTok{\textless{}}\NormalTok{ HEIGHT}\OperatorTok{;}\NormalTok{ j}\OperatorTok{++}\NormalTok{) \{}
                \KeywordTok{const}\NormalTok{ pos }\OperatorTok{=}\NormalTok{ j }\OperatorTok{*}\NormalTok{ WIDTH }\OperatorTok{+}\NormalTok{ i}\OperatorTok{;}
                \ControlFlowTok{if}\NormalTok{ (solution2[j] }\OperatorTok{\textless{}} \DecValTok{0}\NormalTok{) \{}
\NormalTok{                    reachedNegativeValue }\OperatorTok{=} \KeywordTok{true}\OperatorTok{;}
\NormalTok{                \}}
\NormalTok{                currentConcentrationData[pos] }\OperatorTok{=}\NormalTok{ solution2[j]}\OperatorTok{;}
\NormalTok{            \}}
\NormalTok{        \}}
\NormalTok{    \}}

    \ControlFlowTok{if}\NormalTok{ (reachedNegativeValue }\OperatorTok{\&\&} \OperatorTok{!}\NormalTok{allowNegativeValues) \{}
        \BuiltInTok{console}\OperatorTok{.}\FunctionTok{warn}\NormalTok{(}\StringTok{"Concentration went negative at ADI"}\NormalTok{)}\OperatorTok{;}
        \ControlFlowTok{return} \KeywordTok{null}\OperatorTok{;}
\NormalTok{    \}}
    \ControlFlowTok{return}\NormalTok{ currentConcentrationData}\OperatorTok{;}
\NormalTok{\}}\OperatorTok{;}
\end{Highlighting}
\end{Shaded}

\section{Assembling the module}\label{assembling-the-module}

The final code block tells Entangled to write the file
\texttt{src/adi.js} by concatenating all the pieces defined above:

\begin{Shaded}
\begin{Highlighting}[]
\OperatorTok{\textless{}\textless{}}\NormalTok{adi}\OperatorTok{{-}}\NormalTok{imports}\OperatorTok{\textgreater{}\textgreater{}}

\OperatorTok{\textless{}\textless{}}\NormalTok{adi}\OperatorTok{{-}}\NormalTok{variables}\OperatorTok{\textgreater{}\textgreater{}}

\OperatorTok{\textless{}\textless{}}\NormalTok{adi}\OperatorTok{{-}}\KeywordTok{set}\OperatorTok{{-}}\NormalTok{properties}\OperatorTok{\textgreater{}\textgreater{}}

\OperatorTok{\textless{}\textless{}}\NormalTok{adi}\OperatorTok{{-}}\NormalTok{update}\OperatorTok{{-}}\NormalTok{sinks}\OperatorTok{{-}}\NormalTok{sources}\OperatorTok{\textgreater{}\textgreater{}}

\OperatorTok{\textless{}\textless{}}\NormalTok{adi}\OperatorTok{{-}}\NormalTok{update}\OperatorTok{{-}}\NormalTok{diagonal}\OperatorTok{{-}}\NormalTok{x}\OperatorTok{\textgreater{}\textgreater{}}

\OperatorTok{\textless{}\textless{}}\NormalTok{adi}\OperatorTok{{-}}\NormalTok{update}\OperatorTok{{-}}\NormalTok{diagonal}\OperatorTok{{-}}\NormalTok{y}\OperatorTok{\textgreater{}\textgreater{}}

\OperatorTok{\textless{}\textless{}}\NormalTok{adi}\OperatorTok{{-}}\NormalTok{main}\OperatorTok{{-}}\KeywordTok{function}\OperatorTok{\textgreater{}\textgreater{}}
\end{Highlighting}
\end{Shaded}

\bookmarksetup{startatroot}

\chapter{Explaining the Crank-Nicolson
Method}\label{explaining-the-crank-nicolson-method}

The Crank‑Nicolson method is a finite‑difference scheme for solving the
1‑D diffusion equation with optional decay and source terms. It is
implicit in time and second‑order accurate in both space and time. The
method leads to a tridiagonal system that is solved efficiently with the
Thomas algorithm.

The module keeps all working arrays and coefficients at module scope so
that they are allocated only once and reused across many calls to
\texttt{CrankNicolson}.

\section{Imports}\label{imports}

We import the Thomas algorithm solver that will be used to solve the
tridiagonal system at each time step:

\phantomsection\label{cn-imports}%
\begin{Shaded}
\begin{Highlighting}[]
\ImportTok{import}\NormalTok{ \{ thomasAlgorithm \} }\ImportTok{from} \StringTok{"./thomasAlgorithm.js"}\OperatorTok{;}
\end{Highlighting}
\end{Shaded}

\section{Module‑level state}\label{modulelevel-state}

The following variables hold the grid length, the tridiagonal
coefficients for the left‑hand side matrix, the work arrays for the
Thomas algorithm, the solution arrays \texttt{u} and \texttt{uNext}, and
the pre‑computed constants \texttt{lambda}, \texttt{beta},
\texttt{halfLambda}, \texttt{centerCoeff}, and \texttt{dt}.

\phantomsection\label{cn-variables}%
\begin{Shaded}
\begin{Highlighting}[]
\CommentTok{// Module{-}level variables for Crank{-}Nicolson}
\KeywordTok{let}\NormalTok{ LENGTH}\OperatorTok{;}
\KeywordTok{let}\NormalTok{ lowerDiagonal}\OperatorTok{,}\NormalTok{ mainDiagonal}\OperatorTok{,}\NormalTok{ upperDiagonal}\OperatorTok{;}
\KeywordTok{let}\NormalTok{ modifiedUpper}\OperatorTok{,}\NormalTok{ modifiedRHS}\OperatorTok{,}\NormalTok{ solution}\OperatorTok{;}
\KeywordTok{let}\NormalTok{ u}\OperatorTok{,}\NormalTok{ uNext}\OperatorTok{;}
\KeywordTok{let}\NormalTok{ lambda}\OperatorTok{,}\NormalTok{ beta}\OperatorTok{,}\NormalTok{ halfLambda}\OperatorTok{,}\NormalTok{ centerCoeff}\OperatorTok{;}
\KeywordTok{let}\NormalTok{ dt}\OperatorTok{;}
\end{Highlighting}
\end{Shaded}

\section{Setting up the Crank‑Nicolson
coefficients}\label{setting-up-the-cranknicolson-coefficients}

\texttt{setCNProperties} must be called once before any simulation. It
stores the grid length and computes the coefficients that define the
tridiagonal system. Neumann (zero‑flux) boundary conditions are
implemented by adjusting the first and last rows of the matrix using
ghost‑cell mirroring.

\phantomsection\label{cn-set-properties}%
\begin{Shaded}
\begin{Highlighting}[]
\CommentTok{/**}
\CommentTok{ * Initialize properties for Crank{-}Nicolson method}
\CommentTok{ * }\AnnotationTok{@param}\CommentTok{ }\CommentVarTok{\{number\}}\CommentTok{ length {-} Number of grid points}
\CommentTok{ * }\AnnotationTok{@param}\CommentTok{ }\CommentVarTok{\{number\}}\CommentTok{ diffusionCoefficient {-} Diffusion coefficient (alpha)}
\CommentTok{ * }\AnnotationTok{@param}\CommentTok{ }\CommentVarTok{\{number\}}\CommentTok{ deltaX {-} Spatial step size}
\CommentTok{ * }\AnnotationTok{@param}\CommentTok{ }\CommentVarTok{\{number\}}\CommentTok{ deltaT {-} Time step size}
\CommentTok{ * }\AnnotationTok{@param}\CommentTok{ }\CommentVarTok{\{number\}}\CommentTok{ decayRate {-} Decay rate constant (k)}
\CommentTok{ */}
\ImportTok{export} \KeywordTok{const}\NormalTok{ setCNProperties }\OperatorTok{=}\NormalTok{ (}
\NormalTok{    length}\OperatorTok{,}
\NormalTok{    diffusionCoefficient}\OperatorTok{,}
\NormalTok{    deltaX}\OperatorTok{,}
\NormalTok{    deltaT}\OperatorTok{,}
\NormalTok{    decayRate }\OperatorTok{=} \DecValTok{0}
\NormalTok{) }\KeywordTok{=\textgreater{}}\NormalTok{ \{}
\NormalTok{    LENGTH }\OperatorTok{=}\NormalTok{ length}\OperatorTok{;}
\NormalTok{    dt }\OperatorTok{=}\NormalTok{ deltaT}\OperatorTok{;}

    \CommentTok{// Calculate coefficients}
\NormalTok{    lambda }\OperatorTok{=}\NormalTok{ (diffusionCoefficient }\OperatorTok{*}\NormalTok{ deltaT) }\OperatorTok{/}\NormalTok{ (deltaX }\OperatorTok{*}\NormalTok{ deltaX)}\OperatorTok{;}
\NormalTok{    beta }\OperatorTok{=}\NormalTok{ (decayRate }\OperatorTok{*}\NormalTok{ deltaT) }\OperatorTok{/} \FloatTok{2.0}\OperatorTok{;}
\NormalTok{    halfLambda }\OperatorTok{=} \FloatTok{0.5} \OperatorTok{*}\NormalTok{ lambda}\OperatorTok{;}
\NormalTok{    centerCoeff }\OperatorTok{=} \FloatTok{1.0} \OperatorTok{+}\NormalTok{ lambda }\OperatorTok{+}\NormalTok{ beta}\OperatorTok{;}

    \CommentTok{// Tridiagonal coefficients for A (left{-}hand side)}
\NormalTok{    lowerDiagonal }\OperatorTok{=} \KeywordTok{new} \BuiltInTok{Float64Array}\NormalTok{(LENGTH)}\OperatorTok{;}
\NormalTok{    mainDiagonal }\OperatorTok{=} \KeywordTok{new} \BuiltInTok{Float64Array}\NormalTok{(LENGTH)}\OperatorTok{;}
\NormalTok{    upperDiagonal }\OperatorTok{=} \KeywordTok{new} \BuiltInTok{Float64Array}\NormalTok{(LENGTH)}\OperatorTok{;}

    \CommentTok{// Interior coefficients}
    \ControlFlowTok{for}\NormalTok{ (}\KeywordTok{let}\NormalTok{ i }\OperatorTok{=} \DecValTok{1}\OperatorTok{;}\NormalTok{ i }\OperatorTok{\textless{}}\NormalTok{ LENGTH }\OperatorTok{{-}} \DecValTok{1}\OperatorTok{;}\NormalTok{ i}\OperatorTok{++}\NormalTok{) \{}
\NormalTok{        lowerDiagonal[i] }\OperatorTok{=} \OperatorTok{{-}}\NormalTok{halfLambda}\OperatorTok{;}
\NormalTok{        mainDiagonal[i] }\OperatorTok{=}\NormalTok{ centerCoeff}\OperatorTok{;}
\NormalTok{        upperDiagonal[i] }\OperatorTok{=} \OperatorTok{{-}}\NormalTok{halfLambda}\OperatorTok{;}
\NormalTok{    \}}

    \CommentTok{// Neumann (reflective) boundaries using centered differences with ghost cells}
    \CommentTok{// u\_\{{-}1\} = u\_1 and u\_n = u\_\{n{-}2\} for zero flux}
    \CommentTok{// Left boundary i = 0: stencil becomes (1+λ+β)u\_0 {-} λu\_1}
\NormalTok{    lowerDiagonal[}\DecValTok{0}\NormalTok{] }\OperatorTok{=} \FloatTok{0.0}\OperatorTok{;}
\NormalTok{    mainDiagonal[}\DecValTok{0}\NormalTok{] }\OperatorTok{=}\NormalTok{ centerCoeff}\OperatorTok{;}
\NormalTok{    upperDiagonal[}\DecValTok{0}\NormalTok{] }\OperatorTok{=} \OperatorTok{{-}}\NormalTok{lambda}\OperatorTok{;}

    \CommentTok{// Right boundary i = LENGTH{-}1: stencil becomes {-}λu\_\{n{-}2\} + (1+λ+β)u\_\{n{-}1\}}
\NormalTok{    lowerDiagonal[LENGTH }\OperatorTok{{-}} \DecValTok{1}\NormalTok{] }\OperatorTok{=} \OperatorTok{{-}}\NormalTok{lambda}\OperatorTok{;}
\NormalTok{    mainDiagonal[LENGTH }\OperatorTok{{-}} \DecValTok{1}\NormalTok{] }\OperatorTok{=}\NormalTok{ centerCoeff}\OperatorTok{;}
\NormalTok{    upperDiagonal[LENGTH }\OperatorTok{{-}} \DecValTok{1}\NormalTok{] }\OperatorTok{=} \FloatTok{0.0}\OperatorTok{;}

    \CommentTok{// Work arrays for Thomas algorithm}
\NormalTok{    modifiedUpper }\OperatorTok{=} \KeywordTok{new} \BuiltInTok{Float64Array}\NormalTok{(LENGTH)}\OperatorTok{;}
\NormalTok{    modifiedRHS }\OperatorTok{=} \KeywordTok{new} \BuiltInTok{Float64Array}\NormalTok{(LENGTH)}\OperatorTok{;}
\NormalTok{    solution }\OperatorTok{=} \KeywordTok{new} \BuiltInTok{Float64Array}\NormalTok{(LENGTH)}\OperatorTok{;}

    \CommentTok{// Solution arrays}
\NormalTok{    u }\OperatorTok{=} \KeywordTok{new} \BuiltInTok{Float64Array}\NormalTok{(LENGTH)}\OperatorTok{;}
\NormalTok{    uNext }\OperatorTok{=} \KeywordTok{new} \BuiltInTok{Float64Array}\NormalTok{(LENGTH)}\OperatorTok{;}
\NormalTok{\}}\OperatorTok{;}
\end{Highlighting}
\end{Shaded}

\section{The time‑stepping loop}\label{the-timestepping-loop}

The main \texttt{CrankNicolson} function advances the concentration
field by \texttt{totalNumberOfIterations} time steps. At each step it
builds the right‑hand side vector using the explicit part of the scheme
(including source terms and Neumann boundary conditions via ghost‑cell
mirroring), solves the tridiagonal system with the Thomas algorithm, and
then swaps the solution arrays for the next iteration.

If a negative concentration appears and \texttt{allowNegativeValues} is
\texttt{false}, the function throws an error.

\phantomsection\label{cn-main-function}%
\begin{Shaded}
\begin{Highlighting}[]
\CommentTok{/**}
\CommentTok{ * Solve 1D diffusion equation using Crank{-}Nicolson method}
\CommentTok{ * }\AnnotationTok{@param}\CommentTok{ }\CommentVarTok{\{Float64Array\}}\CommentTok{ concentrationData {-} Initial concentration values}
\CommentTok{ * }\AnnotationTok{@param}\CommentTok{ }\CommentVarTok{\{Float64Array\}}\CommentTok{ sources {-} Source terms}
\CommentTok{ * }\AnnotationTok{@param}\CommentTok{ }\CommentVarTok{\{number\}}\CommentTok{ totalNumberOfIterations {-} Number of time steps}
\CommentTok{ * }\AnnotationTok{@param}\CommentTok{ }\CommentVarTok{\{boolean\}}\CommentTok{ allowNegativeValues {-} Whether to allow negative concentrations}
\CommentTok{ */}
\ImportTok{export} \KeywordTok{const}\NormalTok{ CrankNicolson }\OperatorTok{=}\NormalTok{ (}
\NormalTok{    concentrationData}\OperatorTok{,}
\NormalTok{    sources}\OperatorTok{,}
\NormalTok{    totalNumberOfIterations}\OperatorTok{,}
\NormalTok{    allowNegativeValues }\OperatorTok{=} \KeywordTok{false}
\NormalTok{) }\KeywordTok{=\textgreater{}}\NormalTok{ \{}
  

    \CommentTok{// Copy initial condition}
    \ControlFlowTok{for}\NormalTok{ (}\KeywordTok{let}\NormalTok{ i }\OperatorTok{=} \DecValTok{0}\OperatorTok{;}\NormalTok{ i }\OperatorTok{\textless{}}\NormalTok{ LENGTH}\OperatorTok{;}\NormalTok{ i}\OperatorTok{++}\NormalTok{) \{}
\NormalTok{        u[i] }\OperatorTok{=}\NormalTok{ concentrationData[i]}\OperatorTok{;}
\NormalTok{    \}}

    \CommentTok{// Time{-}stepping loop}
    \ControlFlowTok{for}\NormalTok{ (}\KeywordTok{let}\NormalTok{ step }\OperatorTok{=} \DecValTok{0}\OperatorTok{;}\NormalTok{ step }\OperatorTok{\textless{}}\NormalTok{ totalNumberOfIterations}\OperatorTok{;}\NormalTok{ step}\OperatorTok{++}\NormalTok{) \{}
        \CommentTok{// Build right{-}hand side d = B*u\^{}n + dt*S, with Neumann BC via mirroring.}
        \KeywordTok{const}\NormalTok{ d }\OperatorTok{=}\NormalTok{ modifiedRHS}\OperatorTok{;}

        \CommentTok{// Left boundary i = 0, using ghost cell u\_\{{-}1\} = u\_1 for zero flux}
        \CommentTok{// B coefficients: (λ/2) u\_\{i{-}1\}\^{}n + (1 {-} λ {-} β) u\_i\^{}n + (λ/2) u\_\{i+1\}\^{}n}
        \CommentTok{// Substituting u\_\{{-}1\} = u\_1: (λ/2) u\_1 + (1 {-} λ {-} β) u\_0 + (λ/2) u\_1}
\NormalTok{        \{}
            \KeywordTok{const}\NormalTok{ i }\OperatorTok{=} \DecValTok{0}\OperatorTok{;}
            \KeywordTok{const}\NormalTok{ u\_im1 }\OperatorTok{=}\NormalTok{ u[}\DecValTok{1}\NormalTok{]}\OperatorTok{;}       \CommentTok{// ghost cell: u\_\{{-}1\} = u\_1}
            \KeywordTok{const}\NormalTok{ u\_i   }\OperatorTok{=}\NormalTok{ u[}\DecValTok{0}\NormalTok{]}\OperatorTok{;}
            \KeywordTok{const}\NormalTok{ u\_ip1 }\OperatorTok{=}\NormalTok{ u[}\DecValTok{1}\NormalTok{]}\OperatorTok{;}

            \KeywordTok{const}\NormalTok{ rhsVal }\OperatorTok{=}
\NormalTok{                halfLambda }\OperatorTok{*}\NormalTok{ u\_im1 }\OperatorTok{+}
\NormalTok{                (}\FloatTok{1.0} \OperatorTok{{-}}\NormalTok{ lambda }\OperatorTok{{-}}\NormalTok{ beta) }\OperatorTok{*}\NormalTok{ u\_i }\OperatorTok{+}
\NormalTok{                halfLambda }\OperatorTok{*}\NormalTok{ u\_ip1 }\OperatorTok{+}
\NormalTok{                dt }\OperatorTok{*}\NormalTok{ sources[i]}\OperatorTok{;}

\NormalTok{            d[i] }\OperatorTok{=}\NormalTok{ rhsVal}\OperatorTok{;}
\NormalTok{        \}}

        \CommentTok{// Interior points 1..LENGTH{-}2}
        \ControlFlowTok{for}\NormalTok{ (}\KeywordTok{let}\NormalTok{ i }\OperatorTok{=} \DecValTok{1}\OperatorTok{;}\NormalTok{ i }\OperatorTok{\textless{}}\NormalTok{ LENGTH }\OperatorTok{{-}} \DecValTok{1}\OperatorTok{;}\NormalTok{ i}\OperatorTok{++}\NormalTok{) \{}
            \KeywordTok{const}\NormalTok{ u\_im1 }\OperatorTok{=}\NormalTok{ u[i }\OperatorTok{{-}} \DecValTok{1}\NormalTok{]}\OperatorTok{;}
            \KeywordTok{const}\NormalTok{ u\_i   }\OperatorTok{=}\NormalTok{ u[i]}\OperatorTok{;}
            \KeywordTok{const}\NormalTok{ u\_ip1 }\OperatorTok{=}\NormalTok{ u[i }\OperatorTok{+} \DecValTok{1}\NormalTok{]}\OperatorTok{;}

            \KeywordTok{const}\NormalTok{ rhsVal }\OperatorTok{=}
\NormalTok{                halfLambda }\OperatorTok{*}\NormalTok{ u\_im1 }\OperatorTok{+}
\NormalTok{                (}\FloatTok{1.0} \OperatorTok{{-}}\NormalTok{ lambda }\OperatorTok{{-}}\NormalTok{ beta) }\OperatorTok{*}\NormalTok{ u\_i }\OperatorTok{+}
\NormalTok{                halfLambda }\OperatorTok{*}\NormalTok{ u\_ip1 }\OperatorTok{+}
\NormalTok{                dt }\OperatorTok{*}\NormalTok{ sources[i]}\OperatorTok{;}

\NormalTok{            d[i] }\OperatorTok{=}\NormalTok{ rhsVal}\OperatorTok{;}
\NormalTok{        \}}

        \CommentTok{// Right boundary i = LENGTH{-}1, using ghost cell u\_n = u\_\{n{-}2\} for zero flux}
        \CommentTok{// Substituting u\_n = u\_\{n{-}2\}: (λ/2) u\_\{n{-}2\} + (1 {-} λ {-} β) u\_\{n{-}1\} + (λ/2) u\_\{n{-}2\}}
\NormalTok{        \{}
            \KeywordTok{const}\NormalTok{ i }\OperatorTok{=}\NormalTok{ LENGTH }\OperatorTok{{-}} \DecValTok{1}\OperatorTok{;}
            \KeywordTok{const}\NormalTok{ u\_im1 }\OperatorTok{=}\NormalTok{ u[LENGTH }\OperatorTok{{-}} \DecValTok{2}\NormalTok{]}\OperatorTok{;}
            \KeywordTok{const}\NormalTok{ u\_i   }\OperatorTok{=}\NormalTok{ u[LENGTH }\OperatorTok{{-}} \DecValTok{1}\NormalTok{]}\OperatorTok{;}
            \KeywordTok{const}\NormalTok{ u\_ip1 }\OperatorTok{=}\NormalTok{ u[LENGTH }\OperatorTok{{-}} \DecValTok{2}\NormalTok{]}\OperatorTok{;}   \CommentTok{// ghost cell: u\_n = u\_\{n{-}2\}}

            \KeywordTok{const}\NormalTok{ rhsVal }\OperatorTok{=}
\NormalTok{                halfLambda }\OperatorTok{*}\NormalTok{ u\_im1 }\OperatorTok{+}
\NormalTok{                (}\FloatTok{1.0} \OperatorTok{{-}}\NormalTok{ lambda }\OperatorTok{{-}}\NormalTok{ beta) }\OperatorTok{*}\NormalTok{ u\_i }\OperatorTok{+}
\NormalTok{                halfLambda }\OperatorTok{*}\NormalTok{ u\_ip1 }\OperatorTok{+}
\NormalTok{                dt }\OperatorTok{*}\NormalTok{ sources[i]}\OperatorTok{;}

\NormalTok{            d[i] }\OperatorTok{=}\NormalTok{ rhsVal}\OperatorTok{;}
\NormalTok{        \}}

        \CommentTok{// Solve A u\^{}\{n+1\} = d}
        \FunctionTok{thomasAlgorithm}\NormalTok{(}
\NormalTok{            lowerDiagonal}\OperatorTok{,}
\NormalTok{            mainDiagonal}\OperatorTok{,}
\NormalTok{            upperDiagonal}\OperatorTok{,}
\NormalTok{            d}\OperatorTok{,}
\NormalTok{            LENGTH}\OperatorTok{,}
\NormalTok{            modifiedUpper}\OperatorTok{,}
\NormalTok{            modifiedRHS}\OperatorTok{,}
\NormalTok{            solution}
\NormalTok{        )}\OperatorTok{;}

        \CommentTok{// Copy solution and optionally check for negativity}
        \KeywordTok{let}\NormalTok{ hasNegative }\OperatorTok{=} \KeywordTok{false}\OperatorTok{;}
        \ControlFlowTok{for}\NormalTok{ (}\KeywordTok{let}\NormalTok{ i }\OperatorTok{=} \DecValTok{0}\OperatorTok{;}\NormalTok{ i }\OperatorTok{\textless{}}\NormalTok{ LENGTH}\OperatorTok{;}\NormalTok{ i}\OperatorTok{++}\NormalTok{) \{}
\NormalTok{            uNext[i] }\OperatorTok{=}\NormalTok{ solution[i]}\OperatorTok{;}
            \ControlFlowTok{if}\NormalTok{ (}\OperatorTok{!}\NormalTok{allowNegativeValues }\OperatorTok{\&\&}\NormalTok{ uNext[i] }\OperatorTok{\textless{}} \DecValTok{0}\NormalTok{) \{}
\NormalTok{                hasNegative }\OperatorTok{=} \KeywordTok{true}\OperatorTok{;}
\NormalTok{            \}}
\NormalTok{        \}}

        \ControlFlowTok{if}\NormalTok{ (hasNegative }\OperatorTok{\&\&} \OperatorTok{!}\NormalTok{allowNegativeValues) \{}
            \ControlFlowTok{throw} \KeywordTok{new} \BuiltInTok{Error}\NormalTok{(}
                \StringTok{"Negative concentrations encountered in CrankNicolson step "} \OperatorTok{+}\NormalTok{ step}
\NormalTok{            )}\OperatorTok{;}
\NormalTok{        \}}

        \CommentTok{// Swap u and uNext without reallocating}
        \KeywordTok{const}\NormalTok{ tmp }\OperatorTok{=}\NormalTok{ u}\OperatorTok{;}
\NormalTok{        u }\OperatorTok{=}\NormalTok{ uNext}\OperatorTok{;}
\NormalTok{        uNext }\OperatorTok{=}\NormalTok{ tmp}\OperatorTok{;}
\NormalTok{    \}}

    \CommentTok{// Copy result back to concentrationData}
    \ControlFlowTok{for}\NormalTok{ (}\KeywordTok{let}\NormalTok{ i }\OperatorTok{=} \DecValTok{0}\OperatorTok{;}\NormalTok{ i }\OperatorTok{\textless{}}\NormalTok{ LENGTH}\OperatorTok{;}\NormalTok{ i}\OperatorTok{++}\NormalTok{) \{}
\NormalTok{        concentrationData[i] }\OperatorTok{=}\NormalTok{ u[i]}\OperatorTok{;}
\NormalTok{    \}}
\NormalTok{\}}\OperatorTok{;}
\end{Highlighting}
\end{Shaded}

\section{Assembling the module}\label{assembling-the-module-1}

The final code block tells Entangled to write the file
\texttt{src/crankNicolson.js} by concatenating all the pieces defined
above.

\begin{Shaded}
\begin{Highlighting}[]
\OperatorTok{\textless{}\textless{}}\NormalTok{cn}\OperatorTok{{-}}\NormalTok{imports}\OperatorTok{\textgreater{}\textgreater{}}

\OperatorTok{\textless{}\textless{}}\NormalTok{cn}\OperatorTok{{-}}\NormalTok{variables}\OperatorTok{\textgreater{}\textgreater{}}

\OperatorTok{\textless{}\textless{}}\NormalTok{cn}\OperatorTok{{-}}\KeywordTok{set}\OperatorTok{{-}}\NormalTok{properties}\OperatorTok{\textgreater{}\textgreater{}}

\OperatorTok{\textless{}\textless{}}\NormalTok{cn}\OperatorTok{{-}}\NormalTok{main}\OperatorTok{{-}}\KeywordTok{function}\OperatorTok{\textgreater{}\textgreater{}}
\end{Highlighting}
\end{Shaded}

\bookmarksetup{startatroot}

\chapter{Analytic Steady‑State
Solution}\label{analytic-steadystate-solution}

For a 2‑D diffusion equation with Neumann (zero‑flux) boundary
conditions, the steady‑state concentration can be expressed as a double
cosine series. Each term in the series is an eigenfunction of the
Laplacian, and the coefficients are obtained by projecting the source
term onto the eigenfunctions.

The module provides a single exported function,
\texttt{analyticSteadyState}, which computes the steady‑state field up
to a given maximum mode number. A helper function,
\texttt{constantSourceTermOptimized}, efficiently computes the
projection integral for a given mode pair \texttt{(n,\ m)}.

\section{Helper: projecting the source
term}\label{helper-projecting-the-source-term}

\texttt{constantSourceTermOptimized} computes the coefficient

{[} Q\_\{nm\} = \frac{4 \epsilon_n \epsilon_m \Delta x^2}{L_x L_y}
\sum\emph{\{\text{active cells}\} s}\{ij\}
\cos\left(\frac{n\pi x_i}{L_x}\right)
\cos\left(\frac{m\pi y_j}{L_y}\right) {]}

where (\epsilon\_n = 0.5) if (n=0) and (1) otherwise. The function
iterates only over cells where the source is non‑zero, using a
pre‑computed list of active indices and pre‑computed cosine arrays.

\phantomsection\label{as-helper}%
\begin{Shaded}
\begin{Highlighting}[]
\KeywordTok{function} \FunctionTok{constantSourceTermOptimized}\NormalTok{(n}\OperatorTok{,}\NormalTok{ m}\OperatorTok{,}\NormalTok{ Lx}\OperatorTok{,}\NormalTok{ Ly}\OperatorTok{,}\NormalTok{ sources}\OperatorTok{,}\NormalTok{ activeSourceIndices}\OperatorTok{,}\NormalTok{ cosX}\OperatorTok{,}\NormalTok{ cosY}\OperatorTok{,}\NormalTok{ WIDTH}\OperatorTok{,}\NormalTok{ deltaX) \{}
    \KeywordTok{const}\NormalTok{ e\_n }\OperatorTok{=}\NormalTok{ n }\OperatorTok{===} \DecValTok{0} \OperatorTok{?} \FloatTok{0.5} \OperatorTok{:} \DecValTok{1}\OperatorTok{;}
    \KeywordTok{const}\NormalTok{ e\_m }\OperatorTok{=}\NormalTok{ m }\OperatorTok{===} \DecValTok{0} \OperatorTok{?} \FloatTok{0.5} \OperatorTok{:} \DecValTok{1}\OperatorTok{;}
    \KeywordTok{const}\NormalTok{ coefficient }\OperatorTok{=}\NormalTok{ (}\DecValTok{4} \OperatorTok{*}\NormalTok{ e\_n }\OperatorTok{*}\NormalTok{ e\_m }\OperatorTok{*}\NormalTok{ deltaX }\OperatorTok{*}\NormalTok{ deltaX) }\OperatorTok{/}\NormalTok{ (Lx }\OperatorTok{*}\NormalTok{ Ly)}\OperatorTok{;}
    \KeywordTok{let}\NormalTok{ sum }\OperatorTok{=} \DecValTok{0}\OperatorTok{;}

    \ControlFlowTok{for}\NormalTok{ (}\KeywordTok{const}\NormalTok{ idx }\KeywordTok{of}\NormalTok{ activeSourceIndices) \{}
        \KeywordTok{const}\NormalTok{ i }\OperatorTok{=}\NormalTok{ idx }\OperatorTok{\%}\NormalTok{ WIDTH}\OperatorTok{;}
        \KeywordTok{const}\NormalTok{ j }\OperatorTok{=} \BuiltInTok{Math}\OperatorTok{.}\FunctionTok{floor}\NormalTok{(idx }\OperatorTok{/}\NormalTok{ WIDTH)}\OperatorTok{;}
\NormalTok{        sum }\OperatorTok{+=}\NormalTok{ sources[idx] }\OperatorTok{*}\NormalTok{ cosX[n][i] }\OperatorTok{*}\NormalTok{ cosY[m][j]}\OperatorTok{;}
\NormalTok{    \}}

    \ControlFlowTok{return}\NormalTok{ coefficient }\OperatorTok{*}\NormalTok{ sum}\OperatorTok{;}
\NormalTok{\}}
\end{Highlighting}
\end{Shaded}

\section{\texorpdfstring{Main function:
\texttt{analyticSteadyState}}{Main function: analyticSteadyState}}\label{main-function-analyticsteadystate}

The function first identifies all cells with a non‑zero source. If there
are none, it returns a zero field immediately. Otherwise it pre‑computes
the cell‑centre coordinates, the cosine values for all modes and all
grid points, and the squared mode numbers.

The steady‑state solution is then built as a double sum over modes:

{[} C(x,y) = \sum\emph{\{m=0\}\^{}\{M\} \sum}\{n=0\}\^{}\{M\}
\frac{Q_{nm}}{D \lambda_{nm} + k} \cos\left(\frac{n\pi x}{L_x}\right)
\cos\left(\frac{m\pi y}{L_y}\right) {]}

where (\lambda\_\{nm\} = \pi\^{}2\left(\frac{n^2}{L_x^2} +
\frac{m^2}{L_y^2}\right)), (D) is the diffusion rate, and (k) is the
decay rate. Modes with negligible amplitude (below (10\^{}\{-15\})) are
skipped to save time.

\phantomsection\label{as-main}%
\begin{Shaded}
\begin{Highlighting}[]
\ImportTok{export} \KeywordTok{const}\NormalTok{ analyticSteadyState }\OperatorTok{=}\NormalTok{ (}
\NormalTok{    WIDTH}\OperatorTok{,}
\NormalTok{    HEIGHT}\OperatorTok{,}
\NormalTok{    DIFFUSION\_RATE}\OperatorTok{,}
\NormalTok{    DECAY\_RATE}\OperatorTok{,}
\NormalTok{    deltaX}\OperatorTok{,} 
\NormalTok{    sources}\OperatorTok{,} 
\NormalTok{    maxMode }
\NormalTok{) }\KeywordTok{=\textgreater{}}\NormalTok{ \{}
    \KeywordTok{const}\NormalTok{ steadyStateConcentration }\OperatorTok{=} \KeywordTok{new} \BuiltInTok{Float64Array}\NormalTok{(WIDTH }\OperatorTok{*}\NormalTok{ HEIGHT)}\OperatorTok{.}\FunctionTok{fill}\NormalTok{(}\DecValTok{0}\NormalTok{)}\OperatorTok{;}

    \CommentTok{// Precompute non{-}zero source locations}
    \KeywordTok{const}\NormalTok{ activeSourceIndices }\OperatorTok{=}\NormalTok{ []}\OperatorTok{;}
    \ControlFlowTok{for}\NormalTok{ (}\KeywordTok{let}\NormalTok{ idx }\OperatorTok{=} \DecValTok{0}\OperatorTok{;}\NormalTok{ idx }\OperatorTok{\textless{}}\NormalTok{ sources}\OperatorTok{.}\AttributeTok{length}\OperatorTok{;}\NormalTok{ idx}\OperatorTok{++}\NormalTok{) \{}
        \ControlFlowTok{if}\NormalTok{ (sources[idx] }\OperatorTok{!==} \DecValTok{0}\NormalTok{) activeSourceIndices}\OperatorTok{.}\FunctionTok{push}\NormalTok{(idx)}\OperatorTok{;}
\NormalTok{    \}}

    \CommentTok{// Early return if no sources}
    \ControlFlowTok{if}\NormalTok{ (activeSourceIndices}\OperatorTok{.}\AttributeTok{length} \OperatorTok{===} \DecValTok{0}\NormalTok{) \{}
        \ControlFlowTok{return}\NormalTok{ steadyStateConcentration}\OperatorTok{;}
\NormalTok{    \}}

    \KeywordTok{const}\NormalTok{ Lx }\OperatorTok{=}\NormalTok{ WIDTH }\OperatorTok{*}\NormalTok{ deltaX}\OperatorTok{;}
    \KeywordTok{const}\NormalTok{ Ly }\OperatorTok{=}\NormalTok{ HEIGHT }\OperatorTok{*}\NormalTok{ deltaX}\OperatorTok{;}
    \KeywordTok{const}\NormalTok{ piSquared }\OperatorTok{=} \BuiltInTok{Math}\OperatorTok{.}\ConstantTok{PI} \OperatorTok{*} \BuiltInTok{Math}\OperatorTok{.}\ConstantTok{PI}\OperatorTok{;}
    \KeywordTok{const}\NormalTok{ LxSquared }\OperatorTok{=}\NormalTok{ Lx }\OperatorTok{*}\NormalTok{ Lx}\OperatorTok{;}
    \KeywordTok{const}\NormalTok{ LySquared }\OperatorTok{=}\NormalTok{ Ly }\OperatorTok{*}\NormalTok{ Ly}\OperatorTok{;}
    \KeywordTok{const}\NormalTok{ invLxSquared }\OperatorTok{=} \DecValTok{1} \OperatorTok{/}\NormalTok{ LxSquared}\OperatorTok{;}
    \KeywordTok{const}\NormalTok{ invLySquared }\OperatorTok{=} \DecValTok{1} \OperatorTok{/}\NormalTok{ LySquared}\OperatorTok{;}

    \CommentTok{// Precompute all x and y coordinates}
    \KeywordTok{const}\NormalTok{ xCoords }\OperatorTok{=} \KeywordTok{new} \BuiltInTok{Float64Array}\NormalTok{(WIDTH)}\OperatorTok{;}
    \KeywordTok{const}\NormalTok{ yCoords }\OperatorTok{=} \KeywordTok{new} \BuiltInTok{Float64Array}\NormalTok{(HEIGHT)}\OperatorTok{;}
    \ControlFlowTok{for}\NormalTok{ (}\KeywordTok{let}\NormalTok{ i }\OperatorTok{=} \DecValTok{0}\OperatorTok{;}\NormalTok{ i }\OperatorTok{\textless{}}\NormalTok{ WIDTH}\OperatorTok{;}\NormalTok{ i}\OperatorTok{++}\NormalTok{) xCoords[i] }\OperatorTok{=}\NormalTok{ (i }\OperatorTok{+} \FloatTok{0.5}\NormalTok{) }\OperatorTok{*}\NormalTok{ deltaX}\OperatorTok{;}
    \ControlFlowTok{for}\NormalTok{ (}\KeywordTok{let}\NormalTok{ j }\OperatorTok{=} \DecValTok{0}\OperatorTok{;}\NormalTok{ j }\OperatorTok{\textless{}}\NormalTok{ HEIGHT}\OperatorTok{;}\NormalTok{ j}\OperatorTok{++}\NormalTok{) yCoords[j] }\OperatorTok{=}\NormalTok{ (j }\OperatorTok{+} \FloatTok{0.5}\NormalTok{) }\OperatorTok{*}\NormalTok{ deltaX}\OperatorTok{;}

    \CommentTok{// Precompute cosine values for all modes and positions}
    \KeywordTok{const}\NormalTok{ cosX }\OperatorTok{=} \BuiltInTok{Array}\NormalTok{(maxMode }\OperatorTok{+} \DecValTok{1}\NormalTok{)}\OperatorTok{;}
    \KeywordTok{const}\NormalTok{ cosY }\OperatorTok{=} \BuiltInTok{Array}\NormalTok{(maxMode }\OperatorTok{+} \DecValTok{1}\NormalTok{)}\OperatorTok{;}

    \ControlFlowTok{for}\NormalTok{ (}\KeywordTok{let}\NormalTok{ n }\OperatorTok{=} \DecValTok{0}\OperatorTok{;}\NormalTok{ n }\OperatorTok{\textless{}=}\NormalTok{ maxMode}\OperatorTok{;}\NormalTok{ n}\OperatorTok{++}\NormalTok{) \{}
        \KeywordTok{const}\NormalTok{ nPi\_Lx }\OperatorTok{=}\NormalTok{ (}\BuiltInTok{Math}\OperatorTok{.}\ConstantTok{PI} \OperatorTok{*}\NormalTok{ n) }\OperatorTok{/}\NormalTok{ Lx}\OperatorTok{;}
\NormalTok{        cosX[n] }\OperatorTok{=} \KeywordTok{new} \BuiltInTok{Float64Array}\NormalTok{(WIDTH)}\OperatorTok{;}
        \ControlFlowTok{for}\NormalTok{ (}\KeywordTok{let}\NormalTok{ i }\OperatorTok{=} \DecValTok{0}\OperatorTok{;}\NormalTok{ i }\OperatorTok{\textless{}}\NormalTok{ WIDTH}\OperatorTok{;}\NormalTok{ i}\OperatorTok{++}\NormalTok{) \{}
\NormalTok{            cosX[n][i] }\OperatorTok{=} \BuiltInTok{Math}\OperatorTok{.}\FunctionTok{cos}\NormalTok{(nPi\_Lx }\OperatorTok{*}\NormalTok{ xCoords[i])}\OperatorTok{;}
\NormalTok{        \}}
\NormalTok{    \}}

    \ControlFlowTok{for}\NormalTok{ (}\KeywordTok{let}\NormalTok{ m }\OperatorTok{=} \DecValTok{0}\OperatorTok{;}\NormalTok{ m }\OperatorTok{\textless{}=}\NormalTok{ maxMode}\OperatorTok{;}\NormalTok{ m}\OperatorTok{++}\NormalTok{) \{}
        \KeywordTok{const}\NormalTok{ mPi\_Ly }\OperatorTok{=}\NormalTok{ (}\BuiltInTok{Math}\OperatorTok{.}\ConstantTok{PI} \OperatorTok{*}\NormalTok{ m) }\OperatorTok{/}\NormalTok{ Ly}\OperatorTok{;}
\NormalTok{        cosY[m] }\OperatorTok{=} \KeywordTok{new} \BuiltInTok{Float64Array}\NormalTok{(HEIGHT)}\OperatorTok{;}
        \ControlFlowTok{for}\NormalTok{ (}\KeywordTok{let}\NormalTok{ j }\OperatorTok{=} \DecValTok{0}\OperatorTok{;}\NormalTok{ j }\OperatorTok{\textless{}}\NormalTok{ HEIGHT}\OperatorTok{;}\NormalTok{ j}\OperatorTok{++}\NormalTok{) \{}
\NormalTok{            cosY[m][j] }\OperatorTok{=} \BuiltInTok{Math}\OperatorTok{.}\FunctionTok{cos}\NormalTok{(mPi\_Ly }\OperatorTok{*}\NormalTok{ yCoords[j])}\OperatorTok{;}
\NormalTok{        \}}
\NormalTok{    \}}

    \CommentTok{// Precompute squared mode numbers}
    \KeywordTok{const}\NormalTok{ nSquared }\OperatorTok{=} \KeywordTok{new} \BuiltInTok{Float64Array}\NormalTok{(maxMode }\OperatorTok{+} \DecValTok{1}\NormalTok{)}\OperatorTok{;}
    \KeywordTok{const}\NormalTok{ mSquared }\OperatorTok{=} \KeywordTok{new} \BuiltInTok{Float64Array}\NormalTok{(maxMode }\OperatorTok{+} \DecValTok{1}\NormalTok{)}\OperatorTok{;}
    \ControlFlowTok{for}\NormalTok{ (}\KeywordTok{let}\NormalTok{ n }\OperatorTok{=} \DecValTok{0}\OperatorTok{;}\NormalTok{ n }\OperatorTok{\textless{}=}\NormalTok{ maxMode}\OperatorTok{;}\NormalTok{ n}\OperatorTok{++}\NormalTok{) nSquared[n] }\OperatorTok{=}\NormalTok{ n }\OperatorTok{*}\NormalTok{ n}\OperatorTok{;}
    \ControlFlowTok{for}\NormalTok{ (}\KeywordTok{let}\NormalTok{ m }\OperatorTok{=} \DecValTok{0}\OperatorTok{;}\NormalTok{ m }\OperatorTok{\textless{}=}\NormalTok{ maxMode}\OperatorTok{;}\NormalTok{ m}\OperatorTok{++}\NormalTok{) mSquared[m] }\OperatorTok{=}\NormalTok{ m }\OperatorTok{*}\NormalTok{ m}\OperatorTok{;}

    \CommentTok{// Compute steady{-}state solution}
    \ControlFlowTok{for}\NormalTok{ (}\KeywordTok{let}\NormalTok{ m }\OperatorTok{=} \DecValTok{0}\OperatorTok{;}\NormalTok{ m }\OperatorTok{\textless{}=}\NormalTok{ maxMode}\OperatorTok{;}\NormalTok{ m}\OperatorTok{++}\NormalTok{) \{}
        \ControlFlowTok{for}\NormalTok{ (}\KeywordTok{let}\NormalTok{ n }\OperatorTok{=} \DecValTok{0}\OperatorTok{;}\NormalTok{ n }\OperatorTok{\textless{}=}\NormalTok{ maxMode}\OperatorTok{;}\NormalTok{ n}\OperatorTok{++}\NormalTok{) \{}
            \KeywordTok{const}\NormalTok{ eigenvalue }\OperatorTok{=}\NormalTok{ piSquared }\OperatorTok{*}\NormalTok{ (nSquared[n] }\OperatorTok{*}\NormalTok{ invLxSquared }\OperatorTok{+}\NormalTok{ mSquared[m] }\OperatorTok{*}\NormalTok{ invLySquared)}\OperatorTok{;}
            \KeywordTok{const}\NormalTok{ K\_mn }\OperatorTok{=}\NormalTok{ DIFFUSION\_RATE }\OperatorTok{*}\NormalTok{ eigenvalue }\OperatorTok{+}\NormalTok{ DECAY\_RATE}\OperatorTok{;}
            \KeywordTok{const}\NormalTok{ Q\_mn }\OperatorTok{=} \FunctionTok{constantSourceTermOptimized}\NormalTok{(}
\NormalTok{                n}\OperatorTok{,}
\NormalTok{                m}\OperatorTok{,}
\NormalTok{                Lx}\OperatorTok{,}
\NormalTok{                Ly}\OperatorTok{,}
\NormalTok{                sources}\OperatorTok{,}
\NormalTok{                activeSourceIndices}\OperatorTok{,}
\NormalTok{                cosX}\OperatorTok{,}
\NormalTok{                cosY}\OperatorTok{,}
\NormalTok{                WIDTH}\OperatorTok{,}
\NormalTok{                deltaX}
\NormalTok{            )}\OperatorTok{;}
            \KeywordTok{const}\NormalTok{ amplitude }\OperatorTok{=}\NormalTok{ Q\_mn }\OperatorTok{/}\NormalTok{ K\_mn}\OperatorTok{;}

            \CommentTok{// Skip modes with negligible contribution}
            \ControlFlowTok{if}\NormalTok{ (}\BuiltInTok{Math}\OperatorTok{.}\FunctionTok{abs}\NormalTok{(amplitude) }\OperatorTok{\textless{}} \FloatTok{1e{-}15}\NormalTok{) }\ControlFlowTok{continue}\OperatorTok{;}

            \CommentTok{// Compute and accumulate eigenfunction values directly}
            \ControlFlowTok{for}\NormalTok{ (}\KeywordTok{let}\NormalTok{ j }\OperatorTok{=} \DecValTok{0}\OperatorTok{;}\NormalTok{ j }\OperatorTok{\textless{}}\NormalTok{ HEIGHT}\OperatorTok{;}\NormalTok{ j}\OperatorTok{++}\NormalTok{) \{}
                \KeywordTok{const}\NormalTok{ cosYval }\OperatorTok{=}\NormalTok{ cosY[m][j]}\OperatorTok{;}
                \ControlFlowTok{for}\NormalTok{ (}\KeywordTok{let}\NormalTok{ i }\OperatorTok{=} \DecValTok{0}\OperatorTok{;}\NormalTok{ i }\OperatorTok{\textless{}}\NormalTok{ WIDTH}\OperatorTok{;}\NormalTok{ i}\OperatorTok{++}\NormalTok{) \{}
\NormalTok{                    steadyStateConcentration[j }\OperatorTok{*}\NormalTok{ WIDTH }\OperatorTok{+}\NormalTok{ i] }\OperatorTok{+=}\NormalTok{ amplitude }\OperatorTok{*}\NormalTok{ cosX[n][i] }\OperatorTok{*}\NormalTok{ cosYval}\OperatorTok{;}
\NormalTok{                \}}
\NormalTok{            \}}
\NormalTok{        \}}
\NormalTok{    \}}

    \ControlFlowTok{return}\NormalTok{ steadyStateConcentration}\OperatorTok{;}
\NormalTok{\}}\OperatorTok{;}
\end{Highlighting}
\end{Shaded}

\section{Assembling the module}\label{assembling-the-module-2}

The final code block tells Entangled to write the file
\texttt{src/analyticSolution.js} by concatenating the pieces defined
above.

\begin{Shaded}
\begin{Highlighting}[]
\OperatorTok{\textless{}\textless{}}\ImportTok{as}\OperatorTok{{-}}\NormalTok{helper}\OperatorTok{\textgreater{}\textgreater{}}

\OperatorTok{\textless{}\textless{}}\ImportTok{as}\OperatorTok{{-}}\NormalTok{main}\OperatorTok{\textgreater{}\textgreater{}}
\end{Highlighting}
\end{Shaded}

\section{Interactive example}\label{interactive-example}

\begin{Shaded}
\begin{Highlighting}[]
\NormalTok{//| echo: false}

\NormalTok{import \{ analyticSteadyState \} from "../src/analyticSolution.js"}
\end{Highlighting}
\end{Shaded}

\begin{Shaded}
\begin{Highlighting}[]
\NormalTok{//| echo: false}

\NormalTok{viewof diffusionRate = Inputs.range([0.01, 2], \{}
\NormalTok{  label: "Diffusion rate",}
\NormalTok{  step: 0.01,}
\NormalTok{  value: 0.2}
\NormalTok{\})}

\NormalTok{viewof decayRate = Inputs.range([0.01, 2], \{}
\NormalTok{  label: "Decay rate",}
\NormalTok{  step: 0.01,}
\NormalTok{  value: 0.1}
\NormalTok{\})}

\NormalTok{viewof maxMode = Inputs.range([1, 40], \{}
\NormalTok{  label: "Maximum mode",}
\NormalTok{  step: 1,}
\NormalTok{  value: 20}
\NormalTok{\})}
\end{Highlighting}
\end{Shaded}

\begin{Shaded}
\begin{Highlighting}[]
\NormalTok{//| label: fig{-}analytic{-}steady{-}state}
\NormalTok{//| fig{-}cap: "Interactive analytic steady{-}state concentration field"}
\NormalTok{//| echo: false}

\NormalTok{WIDTH = 60;}
\NormalTok{HEIGHT = 40;}
\NormalTok{deltaX = 1;}

\NormalTok{// A single source at the grid centre.}
\NormalTok{sources = new Float64Array(WIDTH * HEIGHT);}
\NormalTok{sources[Math.floor(HEIGHT / 2) * WIDTH + Math.floor(WIDTH / 2)] = 1;}

\NormalTok{concentration = analyticSteadyState(}
\NormalTok{  WIDTH,}
\NormalTok{  HEIGHT,}
\NormalTok{  diffusionRate,}
\NormalTok{  decayRate,}
\NormalTok{  deltaX,}
\NormalTok{  sources,}
\NormalTok{  maxMode}
\NormalTok{);}

\NormalTok{cells = Array.from(concentration, (value, index) =\textgreater{} (\{}
\NormalTok{  x: index \% WIDTH,}
\NormalTok{  y: Math.floor(index / WIDTH),}
\NormalTok{  concentration: value}
\NormalTok{\}));}

\NormalTok{Plot.plot(\{}
\NormalTok{  width,}
\NormalTok{  height: Math.round(width * HEIGHT / WIDTH),}
\NormalTok{  aspectRatio: WIDTH / HEIGHT,}
\NormalTok{  x: \{label: "x"\},}
\NormalTok{  y: \{label: "y", reverse: true\},}
\NormalTok{  color: \{}
\NormalTok{    scheme: "viridis",}
\NormalTok{    label: "Concentration"}
\NormalTok{  \},}
\NormalTok{  marks: [}
\NormalTok{    Plot.cell(cells, \{}
\NormalTok{      x: "x",}
\NormalTok{      y: "y",}
\NormalTok{      fill: "concentration",}
\NormalTok{      inset: 0}
\NormalTok{    \})}
\NormalTok{  ]}
\NormalTok{\})}
\end{Highlighting}
\end{Shaded}

\bookmarksetup{startatroot}

\chapter{}\label{section}

\bookmarksetup{startatroot}

\chapter{}\label{section-1}

\bookmarksetup{startatroot}

\chapter{}\label{section-2}

\bookmarksetup{startatroot}

\chapter{Helper functions}\label{helper-functions}

This module provides utility functions for the diffusion package. Rather
than maintaining one monolithic file, we build \texttt{helpers.js} from
several small, named pieces. Each piece is explained in prose
immediately before or after the code, and Entangled weaves them back
together into a single file when you run \texttt{entangled\ tangle}.

The overall shape of the file is laid out first, using noweb references
(the \texttt{\textless{}\textless{}...\textgreater{}\textgreater{}}
markers) as placeholders for content defined later in this document.

\begin{Shaded}
\begin{Highlighting}[]
\OperatorTok{\textless{}\textless{}}\NormalTok{helpers}\OperatorTok{{-}}\NormalTok{random}\OperatorTok{{-}}\NormalTok{generators}\OperatorTok{\textgreater{}\textgreater{}}

\OperatorTok{\textless{}\textless{}}\NormalTok{helpers}\OperatorTok{{-}}\NormalTok{steady}\OperatorTok{{-}}\NormalTok{state}\OperatorTok{\textgreater{}\textgreater{}}

\OperatorTok{\textless{}\textless{}}\NormalTok{helpers}\OperatorTok{{-}}\NormalTok{grid}\OperatorTok{{-}}\NormalTok{conversion}\OperatorTok{\textgreater{}\textgreater{}}

\OperatorTok{\textless{}\textless{}}\NormalTok{helpers}\OperatorTok{{-}}\NormalTok{grid}\OperatorTok{{-}}\NormalTok{difference}\OperatorTok{\textgreater{}\textgreater{}}
\end{Highlighting}
\end{Shaded}

\section{Generating random sources and
sinks}\label{generating-random-sources-and-sinks}

Diffusion simulations often need randomly scattered ``sources'' (points
that inject material) and ``sinks'' (points that remove it). Both are
built the same way: for every cell in a \texttt{width} × \texttt{height}
grid, we roll a random number and compare it against a probability
threshold to decide whether that cell becomes active.

\texttt{createRandomSources} returns a flat \texttt{Float64Array} where
active cells are marked with \texttt{1.0} and inactive cells with
\texttt{0.0}. Cells are stored in row-major order, so the linear index
for grid position \texttt{(i,\ j)} is \texttt{j\ *\ width\ +\ i}.

\phantomsection\label{helpers-random-generators}%
\begin{Shaded}
\begin{Highlighting}[]
\ImportTok{export} \KeywordTok{const}\NormalTok{ createRandomSources }\OperatorTok{=}\NormalTok{ (width}\OperatorTok{,}\NormalTok{ height}\OperatorTok{,}\NormalTok{ probability}\OperatorTok{,}\NormalTok{ rng }\OperatorTok{=} \BuiltInTok{Math}\OperatorTok{.}\FunctionTok{random}\NormalTok{) }\KeywordTok{=\textgreater{}}\NormalTok{ \{}
    \KeywordTok{const}\NormalTok{ sources }\OperatorTok{=} \KeywordTok{new} \BuiltInTok{Float64Array}\NormalTok{(width }\OperatorTok{*}\NormalTok{ height)}\OperatorTok{;}
    \ControlFlowTok{for}\NormalTok{ (}\KeywordTok{let}\NormalTok{ j }\OperatorTok{=} \DecValTok{0}\OperatorTok{;}\NormalTok{ j }\OperatorTok{\textless{}}\NormalTok{ height}\OperatorTok{;}\NormalTok{ j}\OperatorTok{++}\NormalTok{) \{}
        \ControlFlowTok{for}\NormalTok{ (}\KeywordTok{let}\NormalTok{ i }\OperatorTok{=} \DecValTok{0}\OperatorTok{;}\NormalTok{ i }\OperatorTok{\textless{}}\NormalTok{ width}\OperatorTok{;}\NormalTok{ i}\OperatorTok{++}\NormalTok{) \{}
            \KeywordTok{const}\NormalTok{ idx }\OperatorTok{=}\NormalTok{ j }\OperatorTok{*}\NormalTok{ width }\OperatorTok{+}\NormalTok{ i}\OperatorTok{;}
\NormalTok{            sources[idx] }\OperatorTok{=} \FunctionTok{rng}\NormalTok{() }\OperatorTok{\textless{}}\NormalTok{ probability }\OperatorTok{?} \FloatTok{1.0} \OperatorTok{:} \FloatTok{0.0}\OperatorTok{;}
\NormalTok{        \}}
\NormalTok{    \}}
    \ControlFlowTok{return}\NormalTok{ sources}\OperatorTok{;}
\NormalTok{\}}\OperatorTok{;}
\end{Highlighting}
\end{Shaded}

\texttt{createRandomSinks} follows the identical pattern, but active
cells are marked with a small strength value (\texttt{0.01}) instead of
\texttt{1.0}, representing a gentle, continuous removal rate rather than
an on/off source term.

\phantomsection\label{helpers-random-generators}%
\begin{Shaded}
\begin{Highlighting}[]

\ImportTok{export} \KeywordTok{const}\NormalTok{ createRandomSinks }\OperatorTok{=}\NormalTok{ (width}\OperatorTok{,}\NormalTok{ height}\OperatorTok{,}\NormalTok{ probability}\OperatorTok{,}\NormalTok{ rng }\OperatorTok{=} \BuiltInTok{Math}\OperatorTok{.}\FunctionTok{random}\NormalTok{) }\KeywordTok{=\textgreater{}}\NormalTok{ \{}
    \KeywordTok{const}\NormalTok{ sinks }\OperatorTok{=} \KeywordTok{new} \BuiltInTok{Float64Array}\NormalTok{(width }\OperatorTok{*}\NormalTok{ height)}\OperatorTok{;}
    \ControlFlowTok{for}\NormalTok{ (}\KeywordTok{let}\NormalTok{ j }\OperatorTok{=} \DecValTok{0}\OperatorTok{;}\NormalTok{ j }\OperatorTok{\textless{}}\NormalTok{ height}\OperatorTok{;}\NormalTok{ j}\OperatorTok{++}\NormalTok{) \{}
        \ControlFlowTok{for}\NormalTok{ (}\KeywordTok{let}\NormalTok{ i }\OperatorTok{=} \DecValTok{0}\OperatorTok{;}\NormalTok{ i }\OperatorTok{\textless{}}\NormalTok{ width}\OperatorTok{;}\NormalTok{ i}\OperatorTok{++}\NormalTok{) \{}
            \KeywordTok{const}\NormalTok{ idx }\OperatorTok{=}\NormalTok{ j }\OperatorTok{*}\NormalTok{ width }\OperatorTok{+}\NormalTok{ i}\OperatorTok{;}
\NormalTok{            sinks[idx] }\OperatorTok{=} \FunctionTok{rng}\NormalTok{() }\OperatorTok{\textless{}}\NormalTok{ probability }\OperatorTok{?} \FloatTok{0.01} \OperatorTok{:} \FloatTok{0.0}\OperatorTok{;} \CommentTok{// example sink strength}
\NormalTok{        \}}
\NormalTok{    \}}
    \ControlFlowTok{return}\NormalTok{ sinks}\OperatorTok{;}
\NormalTok{\}}\OperatorTok{;}
\end{Highlighting}
\end{Shaded}

Both functions accept an optional \texttt{rng} parameter defaulting to
\texttt{Math.random}, which makes them easy to test deterministically by
injecting a seeded random number generator.

\section{Detecting steady state}\label{detecting-steady-state}

Iterative diffusion solvers need to know when to stop: once the grid
stops changing meaningfully between iterations, the simulation has
converged. \texttt{checkForSteadyState} compares two successive grids
cell by cell and returns \texttt{true} once the largest absolute
difference between them drops below a \texttt{tolerance} threshold
(default \texttt{1e-5}).

\phantomsection\label{helpers-steady-state}%
\begin{Shaded}
\begin{Highlighting}[]
\ImportTok{export} \KeywordTok{const}\NormalTok{ checkForSteadyState }\OperatorTok{=}\NormalTok{ (prev}\OperatorTok{,}\NormalTok{ current}\OperatorTok{,}\NormalTok{ tolerance }\OperatorTok{=} \FloatTok{1e{-}5}\NormalTok{) }\KeywordTok{=\textgreater{}}\NormalTok{ \{}
    \KeywordTok{let}\NormalTok{ maxDiff }\OperatorTok{=} \DecValTok{0}\OperatorTok{;}
    \ControlFlowTok{for}\NormalTok{ (}\KeywordTok{let}\NormalTok{ i }\OperatorTok{=} \DecValTok{0}\OperatorTok{;}\NormalTok{ i }\OperatorTok{\textless{}}\NormalTok{ prev}\OperatorTok{.}\AttributeTok{length}\OperatorTok{;}\NormalTok{ i}\OperatorTok{++}\NormalTok{) \{}
        \KeywordTok{const}\NormalTok{ diff }\OperatorTok{=} \BuiltInTok{Math}\OperatorTok{.}\FunctionTok{abs}\NormalTok{(current[i] }\OperatorTok{{-}}\NormalTok{ prev[i])}\OperatorTok{;}
        \ControlFlowTok{if}\NormalTok{ (diff }\OperatorTok{\textgreater{}}\NormalTok{ maxDiff) \{}
\NormalTok{            maxDiff }\OperatorTok{=}\NormalTok{ diff}\OperatorTok{;}
\NormalTok{        \}}
\NormalTok{    \}}
    \ControlFlowTok{return}\NormalTok{ maxDiff }\OperatorTok{\textless{}}\NormalTok{ tolerance}\OperatorTok{;}
\NormalTok{\}}\OperatorTok{;}
\end{Highlighting}
\end{Shaded}

Using the maximum difference (rather than, say, an average) is a
conservative choice: it guarantees every single cell has settled, not
just the grid as a whole on average.

\section{Converting flat grids to 2D
matrices}\label{converting-flat-grids-to-2d-matrices}

Internally, grids are stored as flat \texttt{Float64Array}s for
performance, but plotting libraries like Plotly expect nested arrays (a
2D matrix). \texttt{convertTo2D} walks the flat array in row-major order
and rebuilds it as an array of rows, undoing the same
\texttt{j\ *\ width\ +\ i} indexing scheme used when the arrays were
created.

\phantomsection\label{helpers-grid-conversion}%
\begin{Shaded}
\begin{Highlighting}[]
\CommentTok{// Convert 1D arrays to 2D matrices for Plotly}
\ImportTok{export} \KeywordTok{const}\NormalTok{ convertTo2D }\OperatorTok{=}\NormalTok{ (array}\OperatorTok{,}\NormalTok{ width}\OperatorTok{,}\NormalTok{ height) }\KeywordTok{=\textgreater{}}\NormalTok{ \{}
    \KeywordTok{const}\NormalTok{ matrix }\OperatorTok{=}\NormalTok{ []}\OperatorTok{;}
    \ControlFlowTok{for}\NormalTok{ (}\KeywordTok{let}\NormalTok{ j }\OperatorTok{=} \DecValTok{0}\OperatorTok{;}\NormalTok{ j }\OperatorTok{\textless{}}\NormalTok{ height}\OperatorTok{;}\NormalTok{ j}\OperatorTok{++}\NormalTok{) \{}
        \KeywordTok{const}\NormalTok{ row }\OperatorTok{=}\NormalTok{ []}\OperatorTok{;}
        \ControlFlowTok{for}\NormalTok{ (}\KeywordTok{let}\NormalTok{ i }\OperatorTok{=} \DecValTok{0}\OperatorTok{;}\NormalTok{ i }\OperatorTok{\textless{}}\NormalTok{ width}\OperatorTok{;}\NormalTok{ i}\OperatorTok{++}\NormalTok{) \{}
\NormalTok{            row}\OperatorTok{.}\FunctionTok{push}\NormalTok{(array[j }\OperatorTok{*}\NormalTok{ width }\OperatorTok{+}\NormalTok{ i])}\OperatorTok{;}
\NormalTok{        \}}
\NormalTok{        matrix}\OperatorTok{.}\FunctionTok{push}\NormalTok{(row)}\OperatorTok{;}
\NormalTok{    \}}
    \ControlFlowTok{return}\NormalTok{ matrix}\OperatorTok{;}
\NormalTok{\}}\OperatorTok{;}
\end{Highlighting}
\end{Shaded}

\section{Comparing two grids}\label{comparing-two-grids}

Finally, \texttt{calculateDifference} computes a cell-by-cell absolute
difference between two grids of equal size. This is useful for
visualizing where two solutions diverge --- for example, comparing a
numerical solution against an analytic one to check accuracy, or
comparing two different solver methods against each other.

\phantomsection\label{helpers-grid-difference}%
\begin{Shaded}
\begin{Highlighting}[]
\CommentTok{// Compute and display difference}
\ImportTok{export} \KeywordTok{const}\NormalTok{ calculateDifference }\OperatorTok{=}\NormalTok{ (grid1}\OperatorTok{,}\NormalTok{ grid2) }\KeywordTok{=\textgreater{}}\NormalTok{ \{}
    \KeywordTok{const}\NormalTok{ difference }\OperatorTok{=} \KeywordTok{new} \BuiltInTok{Float64Array}\NormalTok{(grid1}\OperatorTok{.}\AttributeTok{length}\NormalTok{)}\OperatorTok{;}
    \ControlFlowTok{for}\NormalTok{ (}\KeywordTok{let}\NormalTok{ i }\OperatorTok{=} \DecValTok{0}\OperatorTok{;}\NormalTok{ i }\OperatorTok{\textless{}}\NormalTok{ grid1}\OperatorTok{.}\AttributeTok{length}\OperatorTok{;}\NormalTok{ i}\OperatorTok{++}\NormalTok{) \{}
\NormalTok{        difference[i] }\OperatorTok{=} \BuiltInTok{Math}\OperatorTok{.}\FunctionTok{abs}\NormalTok{(grid1[i] }\OperatorTok{{-}}\NormalTok{ grid2[i])}\OperatorTok{;}
\NormalTok{    \}}
    \ControlFlowTok{return}\NormalTok{ difference}\OperatorTok{;}
\NormalTok{\}}\OperatorTok{;}
\end{Highlighting}
\end{Shaded}

\bookmarksetup{startatroot}

\chapter*{References}\label{references}
\addcontentsline{toc}{chapter}{References}

\markboth{References}{References}

\phantomsection\label{refs}




\end{document}
